\chapter{Design and Methodology}
\label{ch:method}

\section{System Overview}
\label{sec:overview}

\viki\ is an offline pipeline: demonstrations are recorded, processed and
exported in distinct stages, with no requirement that any stage run at
interactive rates. This is a deliberate design choice. Relaxing the real-time
constraint that shapes teleoperation systems permits offline reconstruction
and robot-plan inspection. Physical replay is a planned validation interface,
not an implemented or evaluated stage.

The pipeline comprises six stages, shown in Fig.~\ref{fig:pipeline}.

\begin{enumerate}[leftmargin=2em]
  \item \emph{Capture.} Synchronised RGB-D streams are recorded from a
        calibrated multi-camera rig while a demonstrator performs the task
        bare-handed.
  \item \emph{Perception.} Hand landmarks are detected per view, lifted to
        three dimensions using sensor depth, and fused across views with
        confidence weights to yield a wrist pose in $\SE$ and a gripper state.
  \item \emph{Representation.} The wrist trajectory is expressed both relative
        to the manipulated object and relative to the ChArUco workspace anchor.
  \item \emph{Retargeting.} A robot joint trajectory is obtained by minimising
        a single weighted cost functional subject to kinematic and collision
        constraints.
  \item \emph{Replay validation (stub).} The interface exists, but no physical
        robot execution or replay-based acceptance is implemented.
  \item \emph{Export.} A NumPy trajectory bundle and JSON manifest can be
        written without implying physical validity; LeRobot export is optional
        and subject to additional requirements.
\end{enumerate}

\begin{figure}[htbp]
\centering
\begin{tikzpicture}[
  node distance=4mm,
  box/.style={draw, rounded corners=1pt, align=center, font=\footnotesize,
              minimum height=8mm, text width=22mm, inner sep=2pt},
  arr/.style={-{Latex[length=2mm]}, thick}]
\node[box] (cap) {Capture\\\scriptsize 2$\times$Kinect};
\node[box, right=of cap] (per) {Perception\\\scriptsize landmarks + fusion};
\node[box, right=of per] (rep) {Representation\\\scriptsize object-relative $\SE$};
\node[box, below=8mm of rep] (ret) {Retargeting\\\scriptsize single cost functional};
\node[box, left=of ret] (val) {Replay\\\scriptsize stub};
\node[box, left=of val] (exp) {Export\\\scriptsize NumPy bundle};
\draw[arr] (cap) -- (per);
\draw[arr] (per) -- (rep);
\draw[arr] (rep) -- (ret);
\draw[arr] (ret) -- (val);
\draw[arr] (val) -- (exp);

\end{tikzpicture}
\caption{Stages of the \viki\ pipeline. Capture, perception and representation
are embodiment-agnostic; retargeting is specific to the target manipulator.
Replay is a stub, and export does not certify physical execution.}
\label{fig:pipeline}
\end{figure}

\section{Hardware and Workspace}
\label{sec:hardware}

The local capture setup uses two Azure Kinect DK cameras observing a
tabletop workspace. UR10 with a parallel-jaw gripper is the intended
simulation target, not a physical manipulator mounted in this setup.
Reachability of the unavailable local demonstrations cannot be quantified.
The Azure Kinect time-of-flight depth characteristics are documented by the
vendor \cite{mskinect} and studied independently \cite{tolgyessy}; neither
source validates the current rig's calibration or fused trajectory accuracy.

The two Kinects are placed to preserve substantial overlap in their fields of
view. This is a direct consequence of the finding of Bu and Lee \cite{bulee}
that iterative-closest-point refinement fails between near-opposing views for
lack of point-cloud overlap; near-orthogonal placement retains both
triangulation baseline and overlap.

\section{Calibration}
\label{sec:calib}

\subsection{Intrinsics and extrinsics}

Intrinsic parameters for every colour and depth stream are estimated from
ChArUco board observations. Extrinsic calibration follows the workspace-anchor
principle: rather than expressing cameras relative to one another, each camera
is calibrated against a ChArUco board rigidly fixed to the workspace, defining
a common world frame $W$. For camera $k$ this yields
\begin{equation}
  \bT^{W}_{C_k} \in \SE,
  \label{eq:extrinsic}
\end{equation}
the transform from the camera frame $C_k$ to the workspace frame $W$. Anchoring
to the workspace rather than to a reference camera has two consequences: adding
or moving a camera does not invalidate the calibration of the others, and the
robot base can be registered into the same frame without a separate
camera-to-camera chain.

\subsection{Robot registration}

The manipulator base frame $R$ is registered to $W$ by observing a ChArUco
target held in a known end-effector pose across a set of configurations, giving
$\bT^{W}_{R}$ directly. Where an explicit camera-to-robot transform is required
instead, the Daniilidis dual-quaternion solver is used, following the
robustness evidence of \cite{alihandeye}, with fifteen to twenty poses spanning
a hemisphere.

\section{Synchronisation}
\label{sec:sync}

The two Azure Kinects are hardware-synchronised in master/subordinate
configuration, with subordinate devices started before the master and depth
exposures staggered to avoid mutual laser interference \cite{mskinect}. The
RealSense cannot share that trigger; its synchronisation mode is incompatible
at the electrical level.

Cross-family alignment therefore uses a single host-side monotonic clock. Every
frame, from every device, is stamped on arrival with
\begin{equation}
  \tau = \texttt{time.monotonic\_ns()},
  \label{eq:hostclock}
\end{equation}
and streams are aligned by nearest neighbour in $\tau$. Two points deserve
emphasis. First, per-device hardware timestamps are \emph{not} used for
cross-device arithmetic: distinct device families count their own cycles from
their own epochs, and the vendor documentation states explicitly that equal
timestamps across cameras are not guaranteed \cite{mskinect}. Second, a
monotonic clock rather than wall-clock time is required, since wall-clock time
may be adjusted discontinuously by time synchronisation daemons, which would
corrupt interval computations mid-recording.

A future local evaluation must preserve paired device and host timestamps to
measure alignment jitter and drift. No local synchronisation statistic is
reported here; the HO-Cap benchmark cannot validate the Kinect rig.

\section{Perception and Multi-View Reconstruction}
\label{sec:fusion}

\subsection{Landmark detection and lifting}

Hand landmarks are detected in each colour stream, yielding for view $k$ at
time $t$ a set of image-plane landmarks with associated visibility scores
$v^{k}_{i,t} \in [0,1]$. Each landmark is lifted to three dimensions using the
registered depth map rather than a learned depth estimate, exploiting the
time-of-flight measurement directly, and transformed into the workspace frame
by \eqref{eq:extrinsic}.

\subsection{Multi-view triangulation}

When a landmark is seen by two or more calibrated views, its position is not
averaged from per-view depth lifts but estimated directly from the image
observations, with depth entering as an additional soft residual rather than as
the anchor. Averaging per-view lifts inherits the depth error of every view;
triangulation is constrained by the camera geometry, which is known to
calibration accuracy and does not degrade with range.

For landmark $i$ at time $t$, let $\mathcal{K}_{i,t}$ be the views whose
detector score clears a threshold and whose observation is finite. Every pair
$(a,b) \subset \mathcal{K}_{i,t}$ yields a linear direct-linear-transform
estimate, which is retained only if it lies in front of both cameras and if the
two viewing rays subtend at least a minimum angle $\theta_{\min}$; a pair with a
short effective baseline is ill-conditioned regardless of how confident the
detector is. Each surviving hypothesis is scored against all views in
$\mathcal{K}_{i,t}$ by the number of reprojection inliers, and the best is
refined over its inlier set by
\begin{equation}
  \hat{\bx}_{i,t} \;=\;
  \argmin_{\bx}\;
  \sum_{k \in \mathcal{K}^{\text{in}}_{i,t}}
  \rho\!\left(\norm{\pi_k(\bx) - \bu^{k}_{i,t}}^{2}\right)
  \;+\;
  \lambda_d \sum_{k} \beta^{k}_{i,t}\,
  \rho_d\!\left(\left(z_k(\bx) - d^{k}_{i,t} - \delta_i\right)^{2}\right),
  \label{eq:triangulate}
\end{equation}
where $\pi_k$ is the pinhole projection of view $k$, $\bu^{k}_{i,t}$ the
observed pixel, $z_k$ the depth of $\bx$ in that camera, $d^{k}_{i,t}$ the
measured depth, and $\delta_i$ a per-landmark offset from the skin surface to
the joint centre---depth sees the skin, the joint sits inside the finger. The
depth weight $\beta^{k}_{i,t}$ is attenuated by the local depth spread and set
to zero for invalid samples, so that a single bad silhouette pixel cannot move
an otherwise well-constrained point. Observations are undistorted once, before
the linear estimate, and every projection downstream is a plain pinhole map;
distortion is never re-applied inside the residual.

The estimate carries a transparent geometric quality
\begin{equation}
  q_{i,t} \;=\;
  \underbrace{\frac{\left|\mathcal{K}^{\text{in}}_{i,t}\right|}
                   {\left|\mathcal{K}_{i,t}\right|}}_{\text{agreement}}
  \cdot
  \underbrace{\operatorname{clip}\!\left(
      1 - \frac{\bar{e}_{i,t}}{2 e_{\text{in}}},\,0,\,1\right)}_{\text{reprojection}}
  \cdot
  \underbrace{\operatorname{clip}\!\left(
      \frac{\theta_{i,t}}{\theta_{\text{ref}}},\,0,\,1\right)}_{\text{baseline}},
  \label{eq:quality}
\end{equation}
with $\bar{e}_{i,t}$ the mean reprojection error over inliers, $\theta_{i,t}$
the widest ray angle, and $e_{\text{in}}$ the inlier radius. This is a
transparent heuristic, not a calibrated variance, and it is used as such: it
weights evidence, and is never interpreted as an inverse covariance.

A property of \eqref{eq:quality} deserves explicit statement, because it
governs how much of a recording survives. The same $e_{\text{in}}$ is the
acceptance threshold, the scale of the robust loss $\rho$ in
\eqref{eq:triangulate}, and the denominator of the reprojection factor. The
factor is therefore designed to grade confidence out to $2 e_{\text{in}}$,
while the threshold admits nothing beyond $e_{\text{in}}$: half of the designed
grading range is unreachable by construction. Widening $e_{\text{in}}$ does not
introduce a soft-weighting mechanism, it permits the existing one to operate
over its own range. A consequence is that $q$ is not comparable across settings
of $e_{\text{in}}$: the same residual scores differently under a different
denominator.

\subsection{Single-view fallback}

Where only one view observes a landmark, triangulation is undefined and the
estimate falls back to the depth lift of that view. The historical multi-view
path also combined per-view lifts by confidence-weighted averaging,
\begin{equation}
  \hat{\bx}_{i,t} \;=\;
  \frac{\sum_{k} w^{k}_{i,t}\, \bx^{k}_{i,t}}
       {\sum_{k} w^{k}_{i,t}},
  \qquad
  w^{k}_{i,t} \;=\;
  \underbrace{v^{k}_{i,t}}_{\text{detector}}
  \cdot
  \underbrace{c^{k}_{i,t}}_{\text{sensor}}
  \cdot
  \underbrace{\frac{1}{\left(d^{k}_{i,t}\right)^{2}}}_{\text{range}}
  \cdot
  \underbrace{\max\!\left(0, \cos\theta^{k}_{i,t}\right)}_{\text{incidence}},
  \label{eq:fusion}
\end{equation}
where $v^{k}_{i,t}$ is the detector score, $c^{k}_{i,t}$ the sensor depth
validity at that pixel, $d^{k}_{i,t}$ the range, and $\theta^{k}_{i,t}$ the
incidence angle. The range and incidence factors follow established practice in
geometric depth fusion, where surfaces observed near perpendicular are weighted
above obliquely observed ones because they are measured more accurately
\cite{fusionpatent}; the detector and sensor factors follow confidence-map
re-weighting \cite{zeng}. The product form ensures that a zero in any
factor---an undetected landmark, an invalid depth pixel, a grazing view---
removes that view's contribution entirely rather than attenuating it.

This path is retained for reproducibility of earlier recordings and for
single-camera operation. It is not the path used for the results reported here,
because averaging per-view lifts propagates depth error into the fused point,
whereas \eqref{eq:triangulate} constrains the point by geometry and admits
depth only as evidence.

\subsection{Wrist pose and gripper state}

The wrist pose $\bT^{W}_{H,t} \in \SE$ is obtained by fitting an orthonormal
frame to the fused wrist and knuckle landmarks. The gripper state is the thumb--index
separation, mapped to a normalised opening,
\begin{equation}
  g_t \;=\;
  \operatorname{clip}\!\left(
    \frac{\norm{\hat{\bx}_{\text{thumb},t} - \hat{\bx}_{\text{index},t}}
          - \tau_{\text{closed}}}
         {\tau_{\text{open}} - \tau_{\text{closed}}},\,0,\,1\right),
  \label{eq:gripper}
\end{equation}
temporally filtered to suppress chatter. An earlier formulation thresholded the
same distance into a binary open/closed indicator; it was replaced because a
parallel-jaw gripper is a continuously commandable actuator, and quantising the
one degree of freedom that carries the grasp discards precisely the approach
and release phases that distinguish a successful demonstration from a failed
one. This single scalar remains the entire hand representation beyond the wrist
frame, following the convention established for parallel-jaw targets in
\cite{opentv, ace}.

\section{External Hand-Joint Reference Protocol}
\label{sec:hocap-method}

The main external cohort contains all 28 selected single-hand HO-Cap
pick-and-place sequences (24 right-hand and four left-hand, 29,706 frames)
from the pinned HO-Cap mirror revision recorded in the archived batch
status. All eight RealSense RGB-D views are processed with MediaPipe observations,
calibrated triangulation and the frozen \texttt{hocap-mediapipe-triangulated-v1}
profile. The HoloLens stream is excluded. These sequences are separate from
the local Kinect cube cohort and cannot validate its calibration. Published
semi-automatic MANO 21-joint positions are transformed into the
\texttt{hocap\_tag\_1} frame. Extract and Prepare do not use the reference
hand labels. Predictions are paired by exact frame index and synthetic
capture timestamp, without root alignment, Procrustes fit, temporal search
or reference-guided filtering. MPJPE is the mean Euclidean error over finite
matched joints; coverage is matched joints divided by valid reference joints.
We report PCK@20/50~mm and the per-sequence 95th percentile alongside MPJPE.
The aggregate MPJPE and PCK are weighted by matched joint counts; an
aggregate P95 cannot be recovered from per-sequence quantiles. Incomplete
Prepare stages are excluded.

For each of the four named objects in a sequence, a separate tracker receives
\emph{published instance masks} and depth from the eight views. It does not
predict segmentation. Model confidence must reach 0.35. A single
right-multiplied object-local SE(3) transform is fit at the first confident
frame using the published object pose. Later matched frames are scored by
translation distance and geodesic SO(3) angle in the same world frame;
coverage is tracked frames divided by valid reference frames. Thus these are
\emph{relative tracking-drift} scores conditional on oracle masks and a
reference-assisted first-frame gauge choice, not absolute 6-DoF localization.
No symmetry-equivalent orientation metric or moving-only subset is applied.
The 28-scene selection is a task-and-handedness subset, not an official
held-out HO-Cap benchmark split. Earlier two-sequence, 120-frame two-versus-
eight-view results remain a development pilot: its camera pair was selected
on the same frames and cannot be treated as held-out evidence.

\section{Object-Centric Representation}
\label{sec:objcentric}

Let $\bT^{W}_{O,t}$ denote the pose of the manipulated object in the workspace
frame. Each demonstration is stored in two forms. The workspace-anchored form
is $\bT^{W}_{H,t}$ directly; the object-relative form is
\begin{equation}
  \bT^{O}_{H,t} \;=\; \left(\bT^{W}_{O,t}\right)^{-1} \bT^{W}_{H,t}.
  \label{eq:objrel}
\end{equation}

The object-relative form is what transfers. If the same task is to be performed
with the object at a new pose $\bT^{R}_{O'}$ expressed in the robot frame, the
desired end-effector pose at time $t$ is
\begin{equation}
  \bT^{R,\mathrm{des}}_{E,t} \;=\;
  \bT^{R}_{O'} \; \bT^{O}_{H,t} \; \bT^{H}_{E},
  \label{eq:transfer}
\end{equation}
where $\bT^{H}_{E}$ is a fixed offset relating the human wrist frame to the
robot end-effector frame, determined once per embodiment. Trajectories are
retained densely rather than reduced to keyframes, since the continuous
velocity profile of the demonstration carries task-relevant information---
approach speed, contact timing---that sparse keyframes discard.

Both representations are exported. Storing both costs little and defers to the
downstream consumer a choice that the literature does not settle universally:
object-relative trajectories generalise across object placement
\cite{orion, spot} but depend on object pose estimation, whereas
workspace-anchored trajectories are simpler and adequate within a fixed layout.

\section{Retargeting as a Single Cost Functional}
\label{sec:retarget}

\subsection{Formulation}

Let $\bq_t \in \R^{n}$ be the manipulator configuration at time $t$ and
$f(\bq_t) \in \SE$ the forward kinematics of the end-effector. The task-space
residual with respect to the desired pose \eqref{eq:transfer} is the
logarithmic map of the pose error,
\begin{equation}
  \be_t(\bq_t) \;=\;
  \Log\!\left(
    \left(\bT^{R,\mathrm{des}}_{E,t}\right)^{-1} f(\bq_t)
  \right)^{\vee} \in \R^{6}.
  \label{eq:residual}
\end{equation}

Retargeting is then posed as the single constrained problem
\begin{equation}
\begin{aligned}
  \min_{\bq_{1:T}} \;\; J(\bq_{1:T}) \;=\; \sum_{t=1}^{T} \Big[\;
   & \underbrace{\omega_t \,\rho_\delta\!\left(
       \norm{\be_t(\bq_t)}_{\bm{\Sigma}}\right)}_{\text{data fidelity}}
   \;+\;
     \underbrace{\lambda_v \norm{\bq_t - \bq_{t-1}}^{2}}_{\text{velocity}}
   \\[1mm]
   &+\;
     \underbrace{\lambda_a \norm{\bq_t - 2\bq_{t-1} + \bq_{t-2}}^{2}}
       _{\text{acceleration}}
   \;+\;
     \underbrace{\lambda_r \norm{\bq_t - \bq_{\mathrm{ref}}}^{2}}
       _{\text{posture}}
  \;\Big]
\end{aligned}
\label{eq:cost}
\end{equation}
subject to
\begin{align}
  \bq_{\min} \;\le\; &\bq_t \;\le\; \bq_{\max}, \label{eq:cjoint}\\
  \norm{\bq_t - \bq_{t-1}} \;&\le\; \dot{\bq}_{\max}\,\Delta t,
    \label{eq:cvel}\\
  h_j(\bq_t) \;&\ge\; 0, \quad j = 1,\dots,m, \label{eq:ccoll}
\end{align}
where

\begin{center}
\begin{tabular}{@{\hspace{10mm}}p{28mm}@{}p{85mm}}
$\omega_t$ & per-frame perception confidence weight, defined in \eqref{eq:conf};\\
$\rho_\delta$ & Huber loss with transition parameter $\delta$;\\
$\bm{\Sigma}$ & diagonal matrix normalising translational and rotational residual units;\\
$\lambda_v, \lambda_a, \lambda_r$ & weights on velocity, acceleration and posture regularisation;\\
$h_j$ & collision and self-collision barrier functions.
\end{tabular}
\end{center}

\subsection{Confidence weighting}

The per-frame weight aggregates the geometric qualities \eqref{eq:quality} of
the landmarks that define the palm frame at time $t$,
\begin{equation}
  \omega_t \;=\;
  \left(\frac{1}{|\mathcal{I}|}\sum_{i \in \mathcal{I}} c_{i,t}\right)^{\alpha}
  \cdot \mathbb{1}\!\left[\text{frame valid}\right],
  \qquad
  c_{i,t} =
  \begin{cases}
    q_{i,t}, & \text{landmark triangulated,}\\
    0, & \text{landmark interpolated,}
  \end{cases}
  \label{eq:conf}
\end{equation}
with $\mathcal{I}$ the four landmarks defining the palm frame---wrist and the
index, middle and little metacarpophalangeal joints---and $\alpha$ controlling
how sharply low-confidence frames are discounted.

Equation \eqref{eq:conf} as written admits a frame whose palm was only
partially observed: a mean over four terms is non-zero whenever one of them is,
so a palm with a single triangulated landmark still contributes at roughly a
quarter weight, and the orientation it contributes was fitted partly through
interpolated points. A stricter diagnostic variant zeroes the frame unless every landmark in
$\mathcal{I}$ was triangulated,
\begin{equation}
  \omega^{\text{strict}}_t \;=\;
  \omega_t \cdot
  \mathbb{1}\!\left[\,\forall i \in \mathcal{I}:\; \text{$i$ triangulated}\,\right].
  \label{eq:confstrict}
\end{equation}
This variant is a proposed diagnostic, not a measured improvement in the
available evidence. It trades observation support for stricter provenance;
a retained local cohort is needed to quantify the effect on tracking and
coverage. The construction is analogous to distance-dependent switching
weights in online teleoperation retargeting \cite{opentv}: poorly observed
frames should exert less influence than well-observed frames.

\subsection{Robustification}

The Huber loss $\rho_\delta$ bounds the influence of gross outliers---a
mis-detected hand, a depth dropout that survives the confidence weighting---at
the level of the data term itself, rather than allowing a filtering stage to
smear the outlier across neighbouring frames \cite{dexteleop0}. This is the
mechanism that replaces pre-smoothing: outliers are down-weighted where they
enter, not attenuated after they have already propagated.

\subsection{Why a single functional}

Equation~\eqref{eq:cost} differs from the sequential alternative in one
structurally important respect. The smoothness penalties act on $\bq_t$, in the
space where smoothness is required, and they are traded against data fidelity
by the solver rather than applied blindly beforehand. In the sequential
pipeline, task-space filtering cannot guarantee joint-space smoothness because
the forward kinematics is nonlinear, and it cannot distinguish task-relevant
fast motion from noise because it operates without reference to the data term.
Here the trade-off is explicit: on high-confidence frames $\omega_t$ is large
and the solution tracks the observation; on low-confidence frames the
smoothness terms dominate and the solution interpolates.

Second-order smoothness, expressed by the acceleration term, follows the
$C^2$-continuity argument of DexFlow \cite{dexflow}. A trajectory that is
first-order smooth may still exhibit curvature discontinuities that appear as
jerk on the physical manipulator and as label noise to a policy trained on the
data.

\subsection{Solution}

The problem is solved with Pink on Pinocchio \cite{pink}. Each term of
\eqref{eq:cost} maps onto a weighted task in the quadratic program---a frame
task for the data term, posture and velocity regularisation tasks for the
smoothness terms---while \eqref{eq:cjoint}--\eqref{eq:ccoll} enter as
inequality constraints, with collision avoidance expressed through control
barrier functions. Levenberg--Marquardt damping stabilises the solution near
singularities.

Because the setting is offline, the problem is solved over the whole trajectory
rather than frame by frame, so the acceleration term is well defined at every
interior timestep and the solution is not constrained by the causality
requirement that an online solver faces. The trajectory is initialised from the
previous solution warm-started along $t$, and boundary conditions at $t=1,2$
are handled by replicating the initial configuration.

\subsection{Interpolation}

Where gaps remain after fusion---landmarks lost across all views---the
trajectory is completed by cubic spline interpolation over $\SE$ prior to
optimisation. Dynamic filters that impose a motion model, such as Kalman or
extended Kalman filters, are deliberately not applied: they introduce a prior
about the dynamics that is subsequently re-estimated by the optimiser, and the
interaction between the two priors is difficult to characterise. The optimiser
already imposes the required regularity through
$\lambda_v$ and $\lambda_a$.

\section{Replay Validation}
\label{sec:replay}

Physical replay remains an interface stub: no robot execution, measured
proprioception, contact outcome, or controller tracking residual is available.
Kinematic plans must not be presented as replay-validated demonstrations.
Simulator or physical validation would require recorded execution, contact
checks, explicit success criteria, and rejection of failed episodes; these
results are not part of the present evidence.

\section{Dataset Export}
\label{sec:export}

The implemented default export writes a self-contained NumPy trajectory bundle
and a JSON manifest. It does not imply physical replay or outcome screening.
An optional LeRobot-format export has additional eligibility requirements and
dependencies; neither format makes an unvalidated plan a physical
demonstration.

\section{Assumptions and Simplifications}
\label{sec:assumptions}

The following assumptions delimit the scope of the work and are revisited in
Section~\ref{sec:implications}.

\begin{enumerate}[leftmargin=2em]
  \item \emph{Parallel-jaw target.} The hand is reduced to a wrist pose and a
        single continuous gripper opening \eqref{eq:gripper}. Multi-fingered
        dexterous targets would require the finger-level representation
        rejected in Section~\ref{sec:lit-hand}.
  \item \emph{The gripper body is not the hand.} The target is the midpoint
        between thumb and index tips, but a parallel-jaw gripper interposes a
        body of its own between the flange and that point, and the whole
        assembly must stay above the workspace plane. For tasks performed
        against the table surface the demonstrated pose can therefore be
        kinematically unreachable even when the arm's reach is ample. The
        size of this gap requires evaluation on retained source episodes.
  \item \emph{Rigid objects.} The object-relative representation
        \eqref{eq:objrel} presupposes a rigid body with a well-defined pose.
        Deformable objects would require a flow-based representation
        \cite{novaflow}.
  \item \emph{Single-arm, quasi-static tasks.} The task class is restricted to
        single-arm tabletop manipulation at moderate speed, where the
        quasi-static assumption underlying purely kinematic retargeting holds.
  \item \emph{Fixed workspace.} The ChArUco anchor is rigidly attached to the
        workspace and the camera rig is not moved within a recording session.
  \item \emph{Kinematic-only transfer.} Forces and contact dynamics are not
        estimated. Tasks whose success depends on force regulation are outside
        the scope.
\end{enumerate}
